The system follows a feasibility-first decomposition that separates three concerns:
time assignment, room assignment, and quality improvement.

\subsection{Decomposition Strategy}
\begin{itemize}[leftmargin=*]
  \item \textbf{Stage 1 (time feasibility):} assign each activity to a valid week/day/slot while satisfying all hard constraints.
  \item \textbf{Stage 2 (rooming):} co-optimize rooms inside CP (\texttt{cp\_rooms}), assign rooms afterward (\texttt{greedy}), or iterate between an exact room subproblem and certificate cuts (\texttt{decomposed}).
  \item \textbf{Stage 3 (quality):} improve a feasible schedule using local search under a soft-penalty objective.
\end{itemize}
This decomposition provides a practical trade-off between strictness and runtime.

\subsection{Adaptive Week-Partitioned Scale Mode}
For large instances, the operational \texttt{partitioned} mode exploits the fact that group, staff, room, and generic-resource conflicts are indexed by teaching week in the current model.
It first runs global structural checks for course totals, staff competency, and block-teaching invariants.
Before partitioning, every required distribution relation is inspected: week-separable aggregate and non-overlap rules are split exactly, fixed week relations are decided statically, and a relation that genuinely couples decisions across weeks causes a fail-fast error, as does a forced repeated-weekly pattern.
Cross-week precedence is accepted only when its fixed week order already satisfies the rule; reverse order is rejected.

Each remaining week is compiled into a compact objective-free CP master.
Up to four week partitions run concurrently by default, each with one CP-SAT search worker to bound memory multiplication.
A valid fast room assignment completes the partition.
If room extraction fails, only that week is rebuilt with the exact certificate-guided room subproblem and the remaining shared wall budget.
Required same-room or different-room relations select the exact inner mode immediately.
All rows are merged only when every partition is feasible, after which the original full instance validator rechecks completeness and every hard rule.
The mode does not report a global soft-objective bound: stability, fairness, and other cross-week quality terms require monolithic optimization or a separately reported improvement phase.

\subsection{Reliability Policy}
To reduce solver timeouts in difficult instances, the runtime applies a bounded retry policy:
\begin{enumerate}[leftmargin=*]
  \item objective-enabled CP attempt (quality-aware),
  \item objective-disabled retry in the same room mode (feasibility-oriented),
  \item optional fallback from strict rooming to greedy rooming.
\end{enumerate}
The operational policy remains available, but research comparisons disable it and analyze primary attempts only.
One-worker configurations are intended to be reproducible for fixed seeds; multi-worker CP-SAT runs are labeled separately.

\subsection{Budgeted Anytime Mode}
For interactive use, total solve time budget $B$ is split into
a feasibility phase $B_f$ and an improvement phase $B_i$:
\begin{align*}
B_f + B_i &= B,\\
B_i &\approx 0.2B\ \text{(capped)},\\
B_f &\approx 0.8B\ \text{(floored)}.
\end{align*}
The first phase prioritizes obtaining a valid schedule quickly; the second phase runs short local-search slices to improve penalty.
Each improvement slice uses a capped deterministic pre-pass (group/week sweep with bounded compound lookahead) followed by guided stochastic neighborhoods with bounded restarts/diversification when the search plateaus.

\subsection{Operational Parameters}
The main runtime parameters are:
\begin{itemize}[leftmargin=*]
  \item room mode (\texttt{cp\_rooms}, \texttt{decomposed}, \texttt{partitioned}, or \texttt{greedy}),
  \item objective on/off,
  \item time limit and worker count,
  \item local-search iteration/time budget.
\end{itemize}
The same solver parameterization is used across batch and interactive runs, supporting consistent interpretation of results.

\subsection{Interactive Administration Layer}
The desktop UI exposes model-consistent administrative edits:
\begin{itemize}[leftmargin=*]
  \item hold-and-move operations across days/slots and also across weeks (subject to hard checks),
  \item immediate conflict diagnostics for candidate targets before committing moves,
  \item real-time recomputation of global soft penalty deltas for candidate placements.
\end{itemize}
Staff week-availability checks are enforced in the same validator path used by solver outputs.

\subsection{Custom Instance Authoring}
The custom generator tab supports:
\begin{itemize}[leftmargin=*]
  \item per-program group/course overrides,
  \item per-course LEC/TUT/LAB count overrides (including lec-only, tut-only, and lab-only patterns inferred directly from counts),
  \item staff day/week availability and teaching-course mapping,
  \item room capacity policy as either categorical or exact numeric mode.
\end{itemize}
Custom generation configurations can be saved locally for reuse and exported/imported as JSON profiles.

\subsection{Generation-Time Solvability Safeguard}
In custom scenarios, users may define room sets that accidentally leave some activities without any eligible room
(because of type, specialization tag, or capacity).
To avoid trivial unsatisfiable instances from such configuration mismatches, generation applies a coverage safeguard:
if an activity has zero eligible rooms, a minimal fallback room of the required class is injected automatically.
This does not remove hard constraints during solving; it only ensures that each generated activity has at least one feasible room candidate.
